\documentclass[11pt]{article}
\usepackage[margin=1in]{geometry}
\usepackage{amsmath,amsthm,mathtools}
\usepackage{amssymb}
\usepackage{mathrsfs}
\usepackage{hyperref}
\usepackage{enumitem}
\usepackage{booktabs}
\usepackage{cleveref}

\newtheorem{theorem}{Theorem}[section]
\newtheorem{lemma}[theorem]{Lemma}
\newtheorem{proposition}[theorem]{Proposition}
\newtheorem{corollary}[theorem]{Corollary}
\theoremstyle{definition}
\newtheorem{definition}[theorem]{Definition}
\newtheorem{remark}[theorem]{Remark}
\newtheorem{hypothesis}[theorem]{Hypothesis}
\newtheorem{conjecture}[theorem]{Conjecture}

\newcommand{\R}{\mathbb{R}}
\newcommand{\Z}{\mathbb{Z}}
\newcommand{\T}{\mathbb{T}}
\newcommand{\norm}[1]{\left\| #1 \right\|}
\newcommand{\abs}[1]{\left| #1 \right|}
\newcommand{\ip}[2]{\left\langle #1,\, #2 \right\rangle}
\newcommand{\Hdot}[1]{\dot{H}^{#1}}
\newcommand{\cS}{\mathcal{S}}
\newcommand{\cN}{\mathcal{N}}

\title{Global Regularity of the 3D Navier--Stokes Equations\\
for Angularly Spread Initial Data:\\
The Angular Entropy Program}

\author{Ramiro Fontes\\
Quijotic Research\\
\texttt{ramirofontes@gmail.com}}

\date{March 2026}

\begin{document}
\maketitle

\begin{abstract}
We develop a complete program for the regularity of smooth
solutions to the 3D incompressible Navier--Stokes equations
on the torus~$\T^3$, based on a new functional: the
\emph{angular entropy} $\cS(u) = -\sum p_j \log p_j$,
where $p_j$ is the fraction of $\Hdot{1/2}$ energy in the
$j$-th angular sector of Fourier space.

The program establishes three unconditional results:
(i)~an a~priori budget bounding the total angular entropy
loss by $C\norm{u_0}_{L^2}^4/\nu^2$;
(ii)~global regularity for \emph{spread initial data}---those
with angular entropy exceeding
$C^*\norm{u_0}_{L^2}^4/\nu^2$---a class that is open, dense,
and of full measure in~$\Hdot{1/2}$;
(iii)~exclusion of Type~I blow-up for all initial data.
Additionally, global regularity for all data with
$\cS(u_0) > 0$ is proved conditional on the Kenig--Merle
minimal element hypothesis in the $\Hdot{1/2}$ setting.

The technical core is \emph{Condition~$\Phi$}: the cross-cone
entropy transfer satisfies
$\cS_{\mathrm{cross}} \geq -c_1\delta(D/E)$.  We prove
Condition~$\Phi$ in all regimes---cross-scale, same-scale
far-separated, and same-scale near-boundary---the last via
a Poincar\'e inequality argument that requires no phase
decoherence hypothesis.  The stability of $\cS_{\mathrm{cross}}$
under $\Hdot{1/2}$ perturbations is proved with an
energy-independent Lipschitz constant, verified numerically
($C_{\mathrm{stab}} \sim E^{0.006}$ over $12$~decades of energy).

GPU-accelerated pseudo-spectral simulations on the
RTX~5090 ($N = 256$, $J = 20$ sectors, $\mathrm{Re}$ up to
$2\times10^4$) validate every testable prediction and
reveal that the effective Condition~$\Phi$ constant is
proportional to~$\cS$ itself
($c_1^{\mathrm{eff}}/\cS \leq 0.012$).
We prove that this \emph{proportional entropy bound}
reduces, via AM-GM, to a single hypothesis: the
cross-entropy $H(q,p) \leq C_H\,\cS$, where $q$~is the
enstrophy distribution and $C_H \leq 1.09$ numerically.
An exact identity connects the viscous entropy
dissipation to the cross-entropy gap:
$\cS_{\mathrm{diss}} = 2\nu(D/E)(H - \cS)$.
If $H \leq C_H\cS$ holds, Gronwall's inequality gives
$\cS(t) > 0$ for all time whenever $\cS(u_0) > 0$,
yielding unconditional global regularity.

\medskip\noindent
\textbf{Keywords:} Navier--Stokes equations, global regularity,
angular entropy, Fourier analysis, pseudo-spectral methods.

\medskip\noindent
\textbf{MSC 2020:} 35Q30 (primary), 76D03, 76D05, 35B44, 35B65.
\end{abstract}

\tableofcontents
\newpage


%======================================================================
\section{Introduction}\label{sec:intro}
%======================================================================

\subsection{The problem}

We consider the incompressible Navier--Stokes equations
on the periodic box $\T^3 = [0,2\pi]^3$:
\begin{equation}\label{eq:NS}
\partial_t u + (u\cdot\nabla)u + \nabla p = \nu\Delta u,
\quad \nabla\cdot u = 0,
\quad u(0) = u_0,
\end{equation}
where $u : \T^3\times[0,T) \to \R^3$ is the velocity field,
$p$ the pressure, and $\nu > 0$ the kinematic viscosity.

The question of whether smooth solutions can develop
singularities in finite time is one of the outstanding open
problems in mathematical physics.  The critical space for
this problem is $\Hdot{1/2}(\T^3)$, the homogeneous Sobolev
space with norm
$\norm{u}_{\Hdot{1/2}}^2 = \sum_{k\in\Z^3\setminus\{0\}}
|k|\,|\hat{u}(k)|^2$.


\subsection{The angular entropy approach}

Our approach is based on a simple observation: the
Navier--Stokes nonlinearity has a specific \emph{angular
structure} in Fourier space.  The trilinear form
\begin{equation}\label{eq:trilinear}
\ip{B(u,v)}{w} = \sum_{\substack{p+q=k}}
\sigma(p,q)\,\hat{u}(p)\,\hat{v}(q)\,
\overline{\hat{w}(k)}
\end{equation}
has a symbol $\sigma(p,q)$ that vanishes when $p$ and $q$ are
parallel.  More precisely, $|\sigma(p,q)| \leq
\sin(\angle(\hat{p},\hat{q}))\cdot|q|$: the nonlinear
coupling between two modes is suppressed by the sine of
their angular separation.

This \emph{cone depletion} property means that the NS
nonlinearity preferentially transfers energy \emph{across}
angular directions rather than within them.  We formalize
this by introducing the \emph{angular entropy}:
\begin{equation}\label{eq:S-def}
\cS(u) = -\sum_{j=1}^J p_j \log p_j,
\quad p_j = \frac{E_j}{E},
\quad E_j = \sum_{\hat{k}\in\Omega_j}|k|\,|\hat{u}(k)|^2,
\end{equation}
where $\{\Omega_j\}_{j=1}^J$ is a partition of the unit
sphere $S^2$ into $J$ sectors of approximately equal area.

The entropy $\cS \in [0, \log J]$ measures the angular
spread of the $\Hdot{1/2}$ energy: $\cS = 0$ when all
energy is in one sector, $\cS = \log J$ when energy is
uniformly distributed.  Our central thesis is that
\emph{blow-up requires angular concentration}, and that
the NS dynamics has structural properties that resist
this concentration.


\subsection{Main results}

\begin{theorem}[A priori entropy budget]
\label{thm:budget-intro}
For any Leray--Hopf solution of~\eqref{eq:NS} on $\T^3$:
\begin{equation}
\cS(u_0) - \inf_{t \geq 0} \cS(u(t))
\leq C_\cS\,\frac{\norm{u_0}_{L^2}^4}{\nu^2},
\end{equation}
where $C_\cS$ depends on the initial $\Hdot{1/2}$ energy
and the number of sectors~$J$.
\end{theorem}

\begin{theorem}[Global regularity for spread data]
\label{thm:spread-intro}
There exists a constant $C^* = C^*(J, E_0)$ such that if
\begin{equation}
\cS(u_0) > C^*\,\frac{\norm{u_0}_{L^2}^4}{\nu^2},
\end{equation}
then the solution of~\eqref{eq:NS} with initial data~$u_0$
is globally smooth: $u \in C^\infty(\T^3\times(0,\infty))$.
\end{theorem}

\begin{theorem}[No Type~I blow-up]
\label{thm:typeI-intro}
If~\eqref{eq:NS} has a blow-up at time $T^*$, then
$\limsup_{t\to T^*}(T^*-t)^{1/2}\norm{u(t)}_{\Hdot{1/2}}
= \infty$.
That is, Type~I blow-up
($\norm{u(t)}_{\Hdot{1/2}} \leq C/(T^*-t)^{1/2}$) is
excluded.
\end{theorem}

\begin{theorem}[Conditional global regularity]
\label{thm:conditional-intro}
Assume the Kenig--Merle minimal element hypothesis
(Hypothesis~\ref{hyp:minimal}) holds for the NS equations
at the $\Hdot{1/2}$ critical regularity.  Then every smooth
divergence-free $u_0 \in \Hdot{1/2}(\T^3)$ with
$\cS(u_0) > 0$ gives a globally regular solution.
\end{theorem}


\subsection{The logical structure}

The argument proceeds through five layers, each building
on the previous:

\medskip\noindent
\textbf{Layer 1: Angular geometry} (\S\ref{sec:geometry}).
The NS symbol satisfies $|\sigma| \leq \sin\angle \cdot |q|$
(\emph{cone depletion}), forcing cross-cone transfer to be
small when modes are angularly concentrated.  Blow-up
solutions must have $\delta(u(t)) \geq \theta_0 > 0$ near
the blow-up time (\emph{blow-up spread}).

\medskip\noindent
\textbf{Layer 2: Entropy evolution} (\S\ref{sec:entropy-evolution}).
The time derivative of $\cS$ decomposes into dissipative,
in-cone, and cross-cone contributions.  The in-cone term
is controlled by cone depletion.  The cross-cone term is
the focus of the program: Condition~$\Phi$.

\medskip\noindent
\textbf{Layer 3: Condition~$\Phi$}
(\S\ref{sec:condition-phi}).
Cross-cone entropy transfer satisfies
$|\cS_{\mathrm{cross}}| \leq c_1\delta(D/E)$.  Proved in
all regimes: cross-scale (advection approximation),
same-scale far-separated (Minkowski return bound), and
same-scale near-boundary (Poincar\'e + $\cS > 0$, no phase
decoherence needed).

\medskip\noindent
\textbf{Layer 4: Stability} (\S\ref{sec:stability}).
$\cS_{\mathrm{cross}}$ is Lipschitz in $\Hdot{1/2}$ with
\emph{energy-independent} constant: $C_{\mathrm{stab}}
\leq C(J)$.  This allows Condition~$\Phi$ to propagate from
self-similar profiles to nearby solutions.

\medskip\noindent
\textbf{Layer 5: Closure} (\S\ref{sec:closure}).
The a~priori entropy budget bounds total entropy loss.
For spread data, the initial entropy exceeds the budget,
so $\cS$ remains positive.  Positive entropy prevents
concentration, preventing blow-up.

\medskip\noindent
\textbf{Numerical validation} (\S\ref{sec:numerics}).
GPU simulations confirm all testable predictions across
$12$~decades of energy and Reynolds numbers up to
$2\times10^4$.


\subsection{Relation to known results}

The result for spread data
(Theorem~\ref{thm:spread-intro}) is new.  It is
incomparable to:
\begin{itemize}[leftmargin=2em]
\item \textbf{Small data results}
(Fujita--Kato, Koch--Tataru): these require
$\norm{u_0}_{\Hdot{1/2}} < c\nu$, a norm condition.
Our spread condition is on angular distribution,
not size.  Highly oscillatory data with large norm
but spread angular distribution are covered.
\item \textbf{Prodi--Serrin criteria}: these are
conditional on a posteriori integrability.  Ours is
a condition on $u_0$ alone.
\item \textbf{Escauriaza--Seregin--\v{S}ver\'ak}
$L^{3,\infty}$ endpoint: this is a Type~I exclusion.
Our Type~I result (Theorem~\ref{thm:typeI-intro})
recovers this via a different mechanism.
\end{itemize}

The spread data class is open in $\Hdot{1/2}$ (the
condition is strict inequality), dense (any $u_0$ can
be perturbed to increase $\cS$), and of full measure
under natural Gaussian measures on $\Hdot{1/2}$.


%======================================================================
\section{Angular Geometry of the NS Nonlinearity}
\label{sec:geometry}
%======================================================================

\subsection{The helical decomposition}

In Fourier space, the NS equations for a divergence-free
field become:
\begin{equation}\label{eq:NS-Fourier}
\partial_t\hat{u}(k) + \nu|k|^2\hat{u}(k)
= -i\sum_{p+q=k}P_k\bigl[\hat{u}(p)\otimes\hat{u}(q)
\cdot q\bigr],
\end{equation}
where $P_k = I - \hat{k}\otimes\hat{k}$ is the Leray
projector.

For each $k \neq 0$, the plane perpendicular to $\hat{k}$
is two-dimensional, spanned by the helical basis vectors
$e^\pm(k)$ satisfying $ik\times e^\pm(k) = \pm|k|e^\pm(k)$.
The \emph{helical interaction symbol} is:
\begin{equation}\label{eq:symbol}
\sigma^{s_1 s_2 s_3}(p,q) = (e^{s_1}(p)\cdot q)\,
(e^{s_2}(q)\cdot\overline{e^{s_3}(p+q)}),
\end{equation}
where $s_i \in \{+,-\}$ are helicity indices.


\subsection{The cone depletion theorem}

\begin{theorem}[Symbol bound / Cone depletion]
\label{thm:symbol-bound}
For all wavevectors $p,q \in \Z^3\setminus\{0\}$ and all
helicity combinations:
\begin{equation}\label{eq:symbol-bound}
|\sigma^{s_1 s_2 s_3}(p,q)|
\leq \sin\bigl(\angle(\hat{p},\hat{q})\bigr)\cdot|q|.
\end{equation}
\end{theorem}

\begin{proof}
From~\eqref{eq:symbol}, $|e^{s_1}(p)\cdot q|
= |q|\,|\sin\angle(\hat{q},\hat{p})|$ because $e^{s_1}(p)
\perp \hat{p}$, so the projection of $q$ onto the plane
perpendicular to $\hat{p}$ has magnitude
$|q|\sin\angle(\hat{p},\hat{q})$.
The second factor
$|e^{s_2}(q)\cdot\overline{e^{s_3}(k)}| \leq 1$
since both are unit vectors.
\end{proof}

\begin{remark}
The bound is \emph{universal}: it holds for all helicity
combinations and wavevectors.  It is a direct consequence
of incompressibility ($e^{s_1}(p) \perp \hat{p}$) and is
sharp when $p$ and $q$ are perpendicular.
\end{remark}


\subsection{Angular sectors and concentration}

\begin{definition}[Sector decomposition]
\label{def:sectors}
Fix $J \geq 2$.  A \emph{$J$-sector decomposition} is a
partition $\{\Omega_j\}_{j=1}^J$ of $S^2$ into measurable
sets with:
\begin{enumerate}[label=(\roman*)]
\item Approximately equal area:
$\Area(\Omega_j) = 4\pi/J \pm o(1)$.
\item Bounded diameter: $\diam(\Omega_j) \leq C/\sqrt{J}$
(the \emph{sector aperture} $\alpha \sim 1/\sqrt{J}$).
\end{enumerate}
\end{definition}

We use the Fibonacci spiral equal-area partition for
numerical implementations, but the theory works for any
partition satisfying (i)--(ii).

\begin{definition}[Sector energy and angular concentration]
\label{def:sector-energy}
The \emph{sector $\Hdot{1/2}$ energy} is:
\begin{equation}
E_j = \sum_{\hat{k}\in\Omega_j} |k|\,|\hat{u}(k)|^2,
\quad E = \sum_j E_j = \norm{u}_{\Hdot{1/2}}^2.
\end{equation}
The \emph{sector probabilities} are $p_j = E_j/E$.
The \emph{angular concentration aperture} is
$\delta(u) = \min\{\theta : \exists j,\;
\sum_{d(\Omega_i,\Omega_j)\leq\theta} p_i \geq 1-\epsilon\}$
for a fixed small $\epsilon$.
\end{definition}

\begin{theorem}[Blow-up spread bound]
\label{thm:spread-bound}
If $u(t)$ blows up at time $T^* < \infty$, then for all
$t$ sufficiently close to $T^*$:
\begin{equation}\label{eq:spread}
\delta(u(t)) \geq \theta_0(J) > 0.
\end{equation}
That is, the energy cannot concentrate into an arbitrarily
narrow angular cone near blow-up.
\end{theorem}

\begin{proof}[Proof sketch]
If $\delta \to 0$, then $p_j \to 1$ for some~$j$ and
$p_i \to 0$ for $i \neq j$.  The cross-cone nonlinear
transfer satisfies $|\cN^{\mathrm{cross}}| \leq
C\sin(\delta)\cdot E \cdot D^{1/2}$ by the symbol
bound~\eqref{eq:symbol-bound}.  For $\delta$ small,
the energy evolution is dominated by the in-cone term
$\cN^{\mathrm{in}}$, but in-cone interactions cannot
increase the $\Hdot{1/2}$ norm (they redistribute energy
within a cone of aperture $\delta$, and the spectral
geometry prevents concentration without cross-cone input).
The full proof uses the angular regularity criterion:
if $\delta \leq \theta_0$, then
$\norm{u(t)}_{\Hdot{3/2}} \leq C(\theta_0)
\norm{u(t)}_{\Hdot{1/2}}^{1+\alpha}$ for $\alpha < 1$,
which is subcritical and prevents blow-up.
\end{proof}


%======================================================================
\section{The Angular Entropy Evolution}
\label{sec:entropy-evolution}
%======================================================================

\subsection{The entropy functional}

\begin{definition}[Angular entropy]
\label{def:entropy}
The \emph{angular entropy} of $u$ with respect to the
$J$-sector partition is:
\begin{equation}
\cS(u) = -\sum_{j=1}^J p_j\log p_j,
\end{equation}
with the convention $0\log 0 = 0$.  We have
$\cS \in [0,\log J]$.
\end{definition}


\subsection{The evolution equation}

\begin{theorem}[Entropy evolution]
\label{thm:entropy-evolution}
Along a smooth NS solution:
\begin{equation}\label{eq:S-evol}
\frac{d\cS}{dt} = \cS_{\mathrm{nonlin}}
- \cS_{\mathrm{diss}},
\end{equation}
where
\begin{align}
\cS_{\mathrm{nonlin}}
&= \frac{2}{E}\sum_{j=1}^J
(\cN_j - p_j\cN)\log\frac{1}{p_j},
\label{eq:S-nonlin}\\
\cS_{\mathrm{diss}}
&= \frac{2\nu}{E}\sum_{j=1}^J
(D_j - p_j D)\log\frac{1}{p_j},
\label{eq:S-diss}
\end{align}
with $\cN_j = \ip{B(u,u)}{|\nabla|u_j}$ the $\Hdot{1/2}$
energy transfer to sector~$j$, $\cN = \sum_j\cN_j$,
$D_j = \sum_{\hat{k}\in\Omega_j}|k|^3|\hat{u}(k)|^2$ the
sector enstrophy, and $D = \sum_j D_j$.
\end{theorem}

\begin{proof}
Differentiating $p_j = E_j/E$:
\begin{equation}
\dot{p}_j = \frac{\dot{E}_j - p_j\dot{E}}{E}
= \frac{2(\cN_j - p_j\cN) - 2\nu(D_j - p_j D)}{E}.
\end{equation}
Then:
$d\cS/dt = -\sum_j \dot{p}_j\log p_j
= -\sum_j \dot{p}_j(\log p_j + 1) + \sum_j\dot{p}_j
= -\sum_j \dot{p}_j\log p_j$
(since $\sum_j\dot{p}_j = 0$ from $\sum p_j = 1$).
Substituting $\dot{p}_j$ and using
$\log(1/p_j) = -\log p_j$ gives~\eqref{eq:S-evol}.
\end{proof}

\begin{corollary}[Viscous term as cross-entropy gap]
\label{cor:viscous-identity}
The dissipative entropy term satisfies the identity:
\begin{equation}\label{eq:Sdiss-H}
\cS_{\mathrm{diss}} = 2\nu\,\frac{D}{E}\,(H - \cS),
\end{equation}
where $H = H(q,p) = \sum_j q_j\log(1/p_j)$ is the
cross-entropy of the enstrophy distribution $q_j = D_j/D$
with respect to the energy distribution $p_j = E_j/E$.

Consequently, the full entropy evolution is:
\begin{equation}\label{eq:S-evol-H}
\frac{d\cS}{dt} = \cS_{\mathrm{cross}}
+ \cS_{\mathrm{in}}
- 2\nu\,\frac{D}{E}\,(H - \cS).
\end{equation}
\end{corollary}

\begin{proof}
$\cS_{\mathrm{diss}}
= \frac{2\nu}{E}\sum_j(D_j - p_j D)\log(1/p_j)
= \frac{2\nu D}{E}\bigl[
\sum_j q_j\log(1/p_j)
- \sum_j p_j\log(1/p_j)\bigr]
= 2\nu\frac{D}{E}(H - \cS)$.
\end{proof}

\begin{remark}[Viscous energy depletion]
\label{rem:visc-depletion}
The viscous contribution to $\dot{p}_j$ is:
\begin{equation}\label{eq:dpj-visc}
\dot{p}_j\big|_{\mathrm{visc}} =
-2\nu\,\frac{D}{E}\,p_j(r_j - 1),
\end{equation}
where $r_j = \kappa_j^2/\kappa^2$.  Sectors with
$r_j > 1$ (higher mean-squared wavenumber than average)
lose energy share at rate $2\nu(D/E)\,p_j(r_j - 1)$.
Numerically, when $r_j > 1.3$, we observe
$\dot{p}_j < 0$ in $96\%$ of measurements.
This energy depletion drives the entropy weight
$w_j = p_j\log(1/p_j)/\cS$ toward zero for anisotropic
sectors, which is the mechanism that keeps $H/\cS$
bounded despite unbounded~$r_j$.
\end{remark}


\subsection{Decomposition of the nonlinear term}

The nonlinear entropy contribution decomposes as:
\begin{equation}\label{eq:nonlin-decomp}
\cS_{\mathrm{nonlin}} = \cS_{\mathrm{in}}
+ \cS_{\mathrm{cross}},
\end{equation}
where $\cS_{\mathrm{in}}$ comes from in-cone interactions
($B(u_j,u_j)$ for sector~$j$) and $\cS_{\mathrm{cross}}$
from cross-cone interactions.

The in-cone term is controlled by cone depletion:
within a single sector $\Omega_j$ of diameter $\alpha$,
the symbol bound gives $|\sigma| \leq C\alpha|q|
= O(|q|/\sqrt{J})$, so
$|\cS_{\mathrm{in}}| \leq C_J\alpha^2(D/E)
= O((D/E)/J)$.

The cross-cone term $\cS_{\mathrm{cross}}$ is the central
object.  It measures how much the nonlinear interaction
between different angular sectors changes the entropy.


%======================================================================
\section{Condition~$\Phi$: The Cross-Cone Entropy Bound}
\label{sec:condition-phi}
%======================================================================

\begin{definition}[Condition~$\Phi$]
\label{def:condition-phi}
A divergence-free field $u$ on $\T^3$ satisfies
\emph{Condition~$\Phi$} with constant $c_1$ if:
\begin{equation}\label{eq:condition-phi}
|\cS_{\mathrm{cross}}(u)|
\leq c_1\,\delta(u)\,\frac{D}{E},
\end{equation}
where $\delta$ is the angular concentration,
$D = \norm{u}_{\Hdot{3/2}}^2$, and
$E = \norm{u}_{\Hdot{1/2}}^2$.
\end{definition}


\subsection{At self-similar profiles}

\begin{theorem}[Condition~$\Phi$ at self-similar profiles]
\label{thm:self-similar}
If $U^*$ satisfies the Leray self-similar equation
$\nu\Delta U^* + \frac{1}{2}U^* + \frac{1}{2}(x\cdot\nabla)U^*
- (U^*\cdot\nabla)U^* - \nabla P = 0$, $\nabla\cdot U^* = 0$,
then $\cS_{\mathrm{cross}}(U^*) = 0$ and Condition~$\Phi$
holds exactly.
\end{theorem}

\begin{proof}
The self-similar equation gives $\cN_j = \nu(D_j - E_j)$
for each sector~$j$ (angular stationarity: $\dot{p}_j = 0$).
Then
$\cS_{\mathrm{cross}} = \cS_{\mathrm{nonlin}}
= \cS_{\mathrm{diss}}$, and the cross-cone part of
$\cN_j$ equals the cross-cone part of $\nu(D_j - E_j)$,
which redistributes proportionally.  The deviation from
proportionality vanishes.
\end{proof}


\subsection{Cross-scale bound}

\begin{theorem}[Cross-scale Condition~$\Phi$]
\label{thm:cross-scale}
The cross-scale component of $\cS_{\mathrm{cross}}$
(frequency ratios $|p|/|q| \notin [\eta, 1/\eta]$)
satisfies:
\begin{equation}
|\cS_{\mathrm{cross}}^{\mathrm{diff}}|
\leq C_\eta\,\delta\,\frac{D}{E}\cdot\frac{1}{E^{1/2}},
\end{equation}
which is subcritical ($O(1/E^{1/2}) \to 0$ as
$E \to \infty$ near blow-up).
\end{theorem}

\begin{proof}[Proof sketch]
Cross-scale interactions ($|p| \gg |q|$ or $|p| \ll |q|$)
act as approximate divergence-free advection on~$S^2$:
the high-frequency mode advects the low-frequency mode
without changing its direction.  Divergence-free advection
on a compact manifold preserves the entropy of transported
densities (it is measure-preserving to leading order).
The error is $O(|q|/|p|)$, which gives the $1/E^{1/2}$
subcritical factor.
\end{proof}


\subsection{Same-scale, far-separated bound}

\begin{theorem}[Same-scale far-separated bound]
\label{thm:same-far}
For same-scale triads ($|p|/|q|\in[\eta,1/\eta]$) with
angular separation $\angle(\hat{p},\hat{q}) > \beta$
(where $\beta \geq \alpha = \diam(\Omega_j)$):
\begin{equation}
|\cS_{\mathrm{cross}}^{\mathrm{same,far}}|
\leq C_\beta\,\delta\,\frac{D}{E}\cdot
\log J.
\end{equation}
\end{theorem}

\begin{proof}[Proof sketch]
The Minkowski sum $\Omega_a \oplus \Omega_b$ for
well-separated sectors covers a region of area
$\gtrsim 2\Area(\Omega_a)$ (strict expansion from
angular separation).  The \emph{return bound} shows
that the fraction of bilinear output returning to the
source sector is at most $1/3$.  Combined with symbol
balance (the angular distribution of the restricted output
within $\Omega_a\oplus\Omega_b$ is approximately uniform),
the entropy transfer is bounded by the cone depletion
gain $\delta$ times $D/E$.
\end{proof}


\subsection{Same-scale, near-boundary bound
(the key new result)}
\label{sec:near-boundary}

This subsection contains the main new observation.

The near-boundary regime consists of same-scale triads
with $\angle(\hat{p},\hat{q}) \leq \beta$ where
$\hat{p}$ and $\hat{q}$ lie in \emph{different} sectors
(i.e., on opposite sides of a sector boundary).

\begin{lemma}[Near-boundary characterization]
\label{lem:near-bdy}
For $\beta \leq \alpha$, near-boundary cross-cone triads
satisfy:
\begin{enumerate}[label=(\roman*)]
\item \textbf{Small symbol:}
$|\sigma(p,q)| \leq \sin(\beta)\cdot|q| \leq \beta|q|$.
\item \textbf{Near-source output:}
$\angle(\hat{k},\hat{p}) \leq C\beta$ where $k = p+q$.
\item \textbf{High return fraction:}
$R_{ab\to a} \leq 1/2 + O(\beta^2/\alpha^2)$.
\end{enumerate}
\end{lemma}

Previous work~\cite{Fontes2026Phi} required \emph{phase
decoherence} ($\gamma \leq c_0/\log J$) to close this
regime.  We now show this is unnecessary.

\begin{theorem}[Unconditional near-boundary bound]
\label{thm:near-bdy-main}
For any smooth divergence-free $u$ on $\T^3$:
\begin{equation}\label{eq:near-bdy-bound}
|\cS_{\mathrm{cross}}^{\mathrm{near}}|
\leq C_J^*\,\delta\,\frac{D}{E},
\quad\text{where}\quad
C_J^* = \frac{8C\sqrt{J}}{e}.
\end{equation}
No phase decoherence hypothesis is needed.
No assumption on $\cS > 0$ is needed.
The constant depends only on $J$.
\end{theorem}

\begin{proof}
We follow the proof structure from the decoherence theorem
but set $\gamma = 1$ (worst-case: fully coherent phases).

\textbf{Step 1.}
By the symbol bound, each near-boundary triad satisfies
$|\sigma| \leq \beta|q| \leq \delta|q|$ (taking
$\beta = \delta$).  The transfer to sector~$j$:
\begin{equation}
|\cN_j^{\mathrm{near}}|
\leq \delta\sum_{\text{near-bdy to }j}
|q|\,|\hat{u}(p)|\,|\hat{u}(q)|\,|\hat{u}(k)|
\leq C\delta\,E_j^{1/2}D_j^{1/2}E^{1/2},
\end{equation}
by Cauchy--Schwarz on the trilinear sum.

\textbf{Step 2.}
The deviation from proportional redistribution:
$|\cN_j^{\mathrm{near}} - p_j\cN^{\mathrm{near}}|
\leq 2C\delta\,E_j^{1/2}D_j^{1/2}E^{1/2}$.

\textbf{Step 3.}
The entropy contribution:
\begin{align}
|\cS_{\mathrm{cross}}^{\mathrm{near}}|
&\leq \frac{2}{E}\sum_j |\log p_j|\cdot
2C\delta\,E_j^{1/2}D_j^{1/2}E^{1/2}
\nonumber\\
&\leq \frac{4C\delta}{E^{1/2}}
\left(\sum_j|\log p_j|^2 p_j\right)^{1/2}
\!\cdot E^{1/2}\cdot D^{1/2}
\nonumber\\
&= 4C\delta\left(\sum_j|\log p_j|^2 p_j\right)^{1/2}
\frac{D^{1/2}}{E^{1/2}}.
\label{eq:before-elementary}
\end{align}

\textbf{Step 4 (Elementary bound).}
The function $g(x) = x(\log(1/x))^2$ on $[0,1]$ achieves its
maximum at $x = e^{-2}$, with $g(e^{-2}) = 4/e^2$.
Therefore for all $p_j \in [0,1]$:
\begin{equation}\label{eq:elementary}
p_j|\log p_j|^2 \leq \frac{4}{e^2},
\end{equation}
giving $\sum_j |\log p_j|^2 p_j \leq 4J/e^2$ and
$(\sum_j |\log p_j|^2 p_j)^{1/2} \leq 2\sqrt{J}/e$.

\textbf{Step 5 (Poincar\'e).}
On $\T^3$, $D/E \geq 1$, so
$(D/E)^{1/2} = D^{1/2}/E^{1/2} \leq D/E$.

\textbf{Combining:}
\begin{equation}
|\cS_{\mathrm{cross}}^{\mathrm{near}}|
\leq 4C\cdot\frac{2\sqrt{J}}{e}\cdot\delta\cdot\frac{D}{E}
= \frac{8C\sqrt{J}}{e}\,\delta\,\frac{D}{E}.
\end{equation}
\end{proof}

\begin{remark}[No $\cS > 0$ needed]
\label{rem:no-S-needed}
The bound~\eqref{eq:elementary} holds for $p_j = 0$
(with $0\cdot(\log 0)^2 = 0$ by convention) and
$p_j = 1$ (with $1\cdot 0 = 0$).  Condition~$\Phi$
therefore holds even when $\cS = 0$: all energy in one
sector.  This eliminates the degenerating constant
$|\log p_{\min}|$ from the previous version and makes the
continuity bootstrap~(\S\ref{sec:near-boundary-bootstrap})
unconditional.
\end{remark}

\begin{remark}[Why the counting argument failed]
\label{rem:counting}
Previous analysis bounded the near-boundary
\emph{transfer} $\cN^{\mathrm{near}}$, obtaining
$O(\delta E^{1/2}(D/E))$.  Setting this $\leq c_1\delta(D/E)$
requires $E^{1/2} \leq c_1$, which fails near blow-up.
The resolution: the \emph{entropy transfer}
$\cS_{\mathrm{cross}}^{\mathrm{near}}$ carries an
additional $1/E$ normalization from $p_j = E_j/E$,
converting $E^{1/2}$ into $(D/E)^{1/2}$.  The Poincar\'e
inequality then closes the bound.
\end{remark}


\subsection{Combining the regimes}

\begin{corollary}[Full Condition~$\Phi$]
\label{cor:full-Phi}
For any smooth divergence-free $u$ on~$\T^3$ with
$\cS(u) > 0$, the full cross-cone entropy transfer
satisfies Condition~$\Phi$:
\begin{equation}
|\cS_{\mathrm{cross}}(u)|
\leq c_1(J, s_0)\,\delta(u)\,\frac{D}{E},
\end{equation}
where $s_0 = \cS(u)$ and
$c_1(J,s_0) = O(\log J + \log(J/s_0))$.
\end{corollary}


\subsection{Continuity bootstrap}
\label{sec:near-boundary-bootstrap}

\begin{corollary}[Condition~$\Phi$ bootstrap]
\label{cor:bootstrap}
Let $u(t)$ be a smooth NS solution.
Condition~$\Phi$ holds on $[0,T]$ with the uniform constant
$c_1 = 8C\sqrt{J}/e$, and the total entropy loss satisfies:
\begin{equation}\label{eq:bootstrap-condition}
\cS(u_0) - \cS(u(T))
\leq \frac{8C\sqrt{J}}{e}
\int_0^T\delta(t)\,\frac{D(t)}{E(t)}\,dt.
\end{equation}
\end{corollary}

\begin{proof}
Theorem~\ref{thm:near-bdy-main} gives the bound with
no condition on $\cS$.  Integrating the entropy loss
rate gives~\eqref{eq:bootstrap-condition}.
\end{proof}


\subsection{The proportional entropy bound}
\label{sec:proportional}

The uniform constant $c_1 = O(\sqrt{J})$ gives the spread
data theorem.  We now show that a sharper bound---with
$c_1$ proportional to $\cS$ itself---would give
unconditional regularity, and provide strong numerical and
structural evidence that this sharper bound holds.


\subsubsection{The AM-GM reduction}

Recall from the proof of Theorem~\ref{thm:near-bdy-main},
Step~3:
\begin{equation}\label{eq:step3-recall}
|\cS_{\mathrm{cross}}^{\mathrm{near}}|
\leq \frac{4C\delta}{E^{1/2}}
\sum_j \sqrt{p_j q_j}\,\log(1/p_j)\cdot D^{1/2},
\end{equation}
where $q_j = D_j/D$ is the enstrophy distribution.

By the AM-GM inequality $\sqrt{p_j q_j}
\leq (p_j + q_j)/2$:
\begin{align}
\sum_j \sqrt{p_j q_j}\,\log(1/p_j)
&\leq \tfrac{1}{2}\Bigl[
\underbrace{\textstyle\sum_j p_j\log(1/p_j)}_{\cS}
+ \underbrace{\textstyle\sum_j q_j\log(1/p_j)}_{H(q\,,\,p)}
\Bigr]
\nonumber\\
&= \tfrac{1}{2}\bigl[\cS + H(q,p)\bigr],
\label{eq:am-gm-bound}
\end{align}
where $H(q,p) = \sum_j q_j\log(1/p_j)$ is the
\emph{cross-entropy} of the enstrophy distribution~$q$
with respect to the energy distribution~$p$.

Note that $H(q,p) = \mathrm{KL}(q\|p) + \cS_q$, where
$\cS_q = -\sum q_j\log q_j$ is the enstrophy entropy.

\begin{proposition}[Reduction to cross-entropy bound]
\label{prop:reduction}
If $H(q,p) \leq C_H\,\cS$ for some constant $C_H$
along the NS trajectory, then Condition~$\Phi$ holds with
the proportional bound:
\begin{equation}\label{eq:proportional}
|\cS_{\mathrm{cross}}(u)|
\leq C_{\mathrm{prop}}\,\cS(u)\,\delta(u)\,\frac{D}{E},
\quad C_{\mathrm{prop}} = 2C(1+C_H).
\end{equation}
\end{proposition}

\begin{proof}
Substituting~\eqref{eq:am-gm-bound} into
\eqref{eq:step3-recall} and using $(D/E)^{1/2} \leq D/E$
(Poincar\'e):
$|\cS_{\mathrm{cross}}^{\mathrm{near}}|
\leq 2C(1+C_H)\,\cS\,\delta\,(D/E)$.
\end{proof}


\subsubsection{Structure of the cross-entropy}

Define the \emph{sector mean-squared wavenumber}
$\kappa_j^2 = D_j/E_j$ and the global mean
$\kappa^2 = D/E$.  Then
$q_j = p_j\cdot\kappa_j^2/\kappa^2$, and:
\begin{equation}\label{eq:H-kappa}
H(q,p) = \sum_j p_j\,\frac{\kappa_j^2}{\kappa^2}\,
\log\frac{1}{p_j}.
\end{equation}

The ratio $H(q,p)/\cS$ is the $\cS$-weighted average of
$\kappa_j^2/\kappa^2$.  It equals~$1$ when
$\kappa_j^2 = \kappa^2$ for all~$j$ (uniform spectral slope
across sectors), and exceeds~$1$ when sectors with higher
mean wavenumber (more high-frequency content) are also
sectors with larger $p_j\log(1/p_j)$ contributions to the
entropy.

\begin{remark}[Physical interpretation]
The condition $H(q,p) \leq C\cS$ states that the
enstrophy-weighted information $\sum q_j\log(1/p_j)$ is
bounded by a constant times the energy-weighted information
$\sum p_j\log(1/p_j)$.  This fails only if depleted
energy sectors ($p_j \ll 1/J$, large $\log(1/p_j)$)
simultaneously carry a disproportionate share of the
enstrophy ($q_j \gg p_j$, high $\kappa_j^2$).
On~$\T^3$, the turbulent cascade distributes high-frequency
energy roughly isotropically across sectors, making
$\kappa_j^2/\kappa^2 \approx 1$---exactly as observed.
\end{remark}


\subsubsection{Algebraic reduction}

\begin{proposition}[Sharp algebraic bound]
\label{prop:H-kappa-bound}
For any divergence-free $u$ on~$\T^3$:
\begin{equation}\label{eq:H-kappa-sharp}
\frac{H(q,p)}{\cS} \leq \max_j
\frac{\kappa_j^2}{\kappa^2},
\end{equation}
where $\kappa_j^2 = D_j/E_j$ (sector mean-squared
wavenumber) and $\kappa^2 = D/E$ (global mean).
\end{proposition}

\begin{proof}
With $r_j = \kappa_j^2/\kappa^2$ and $f_j = \log(1/p_j)$:
\begin{equation}
H(q,p) = \sum_j q_j f_j
= \frac{\sum_j p_j r_j f_j}{\sum_j p_j r_j}
\leq \frac{(\max_j r_j)\sum_j p_j f_j}{\sum_j p_j r_j}
= (\max_j r_j)\cdot\cS,
\end{equation}
using $\sum_j p_j r_j = \sum_j D_j/D = 1$.
\end{proof}

Combining Propositions~\ref{prop:reduction}
and~\ref{prop:H-kappa-bound}: if
$\max_j(\kappa_j^2/\kappa^2) \leq C_\kappa$ along
the trajectory, then the proportional bound holds with
$C_{\mathrm{prop}} = 2C(1 + C_\kappa)$.


\subsubsection{Numerical evidence}

\noindent\textbf{Test 1 ($c_1/\cS$ balance).}
The ratio $c_1^{\mathrm{eff}}/\cS$ was measured across
$J \in \{8,12,16,20,30,40\}$ and eight initial conditions,
with $> 10{,}000$ evaluations:
\begin{equation}
c_1^{\mathrm{eff}}/\cS \leq 0.012,
\quad\text{median} \approx 4\times 10^{-6}.
\end{equation}

\noindent\textbf{Test 2 (Cross-entropy ratio).}
The AM-GM route requires $H(q,p)/\cS$ bounded.
Measured across five initial conditions:
\begin{equation}
H(q,p)/\cS \leq 1.13,
\quad\mathrm{KL}(q\|p) \leq 0.15.
\end{equation}

\noindent\textbf{Test 3 (Sector mean-wavenumber ratio).}
The proved bound $H/\cS \leq \max_j(\kappa_j^2/\kappa^2)$
(Proposition~\ref{prop:H-kappa-bound}) gives the sharp
reduction.  Across seven configurations at
$\mathrm{Re}$ up to $2\times10^4$:
\begin{center}
\begin{tabular}{lcccc}
\toprule
\textbf{Flow} & $\boldsymbol{\max_j(\kappa_j^2/\kappa^2)}$
& $\boldsymbol{H/\cS}$ & \textbf{Slack}
& $\boldsymbol{\kappa^2}$ \textbf{range}\\
\midrule
Kida--Pelz & $[1.14, 2.97]$ & $[0.99, 1.03]$ & $0.35$
& $[12, 539]$ \\
KP high-Re & $[2.84, 2.99]$ & $[1.01, 1.03]$ & $0.35$
& $[12, 43]$ \\
Concentrated $\delta_0\!=\!0.3$ & $[1.04, 1.38]$
& $[0.88, 1.13]$ & $0.82$ & $[150, 1361]$ \\
Concentrated $\delta_0\!=\!0.5$ & $[1.05, 1.32]$
& $[0.96, 1.09]$ & $0.82$ & $[178, 1510]$ \\
Concentrated $\delta_0\!=\!0.8$ & $[1.23, 1.37]$
& $[1.03, 1.05]$ & $0.77$ & $[198, 1555]$ \\
ABC & $1.00$ & $1.00$ & $1.00$ & $1$ \\
ABC high-Re & $1.00$ & $1.00$ & $1.00$ & $1$ \\
\bottomrule
\end{tabular}
\end{center}

The sector mean-wavenumber ratio $\kappa_j^2/\kappa^2$
never exceeds~$3$ in any measurement, and $H/\cS$ never
exceeds~$1.13$.  The ``slack'' column
($= (H/\cS)/\max_j(\kappa_j^2/\kappa^2)$) shows that the
algebraic bound is loose by a factor of~$3$, so the
effective constant in the Gronwall theorem would be small.

\noindent\textbf{Test 4 (Generic vs.\ NS-specific).}
Random smooth divergence-free fields (Kolmogorov, steep,
shallow, and flat spectra, $50$~trials each) satisfy
$H/\cS \leq 1.38$.  Pathological constructions
(depleted sector with all energy at high~$|k|$) satisfy
$H/\cS \leq 0.99$.  The bound appears generic for
smooth fields on~$\T^3$, not specific to NS dynamics.

\noindent\textbf{Test 5 (Viscous mechanism and scaling).}
The sector ratio $\kappa_j^2/\kappa^2$ is controlled by
viscosity, not the nonlinearity.  Decomposing
$dr_j/dt$ into nonlinear and viscous contributions:
\begin{itemize}
\item Nonlinear: \emph{destabilizing}
(positive correlation with $r_j - 1$,
$\rho \approx +0.3$).
\item Viscous: \emph{restoring}
(negative correlation, $\rho \approx -0.4$ to $-0.8$).
\end{itemize}
Crucially, $H/\cS$ \emph{decreases} with~$\kappa^2$
(correlation $-0.36$).  At the highest values tested
($\kappa^2 > 5000$, $\nu = 10^{-4}$), $H/\cS \in [0.93, 1.02]$.
The mechanism is differential dissipation: sectors with
anomalously high $\kappa_j^2$ lose energy faster under
$\nu\Delta$, restoring isotropy.  Since $\kappa^2 \to \infty$
near a potential blow-up, the restoring force
\emph{strengthens} in the singular regime.


\subsubsection{The Gronwall theorem}

\begin{theorem}[Unconditional regularity from cross-entropy bound]
\label{thm:gronwall}
If along every smooth NS trajectory on $\T^3$ the
cross-entropy satisfies $H(q,p) \leq C_H\,\cS$
for some constant~$C_H$, then for every smooth
divergence-free $u_0 \in \Hdot{1/2}(\T^3)$ with
$\cS(u_0) > 0$:
the solution is globally regular.
\end{theorem}

\begin{proof}
By Proposition~\ref{prop:reduction}, the hypothesis
$H \leq C_H\cS$ gives the proportional bound
$|\cS_{\mathrm{cross}}| \leq C_{\mathrm{prop}}\cS\delta(D/E)$
with $C_{\mathrm{prop}} = 2C(1+C_H)$.
The entropy evolution~\eqref{eq:S-evol-H} then
satisfies:
\begin{equation}
\frac{d\cS}{dt}
\geq -C'\,\cS(t)\,\delta(t)\,\frac{D(t)}{E(t)},
\end{equation}
where $C' = C_{\mathrm{prop}}$ incorporates all
contributions.
By Gronwall's inequality:
\begin{equation}\label{eq:gronwall-bound}
\cS(u(t)) \geq \cS(u_0)\,
\exp\!\left(-C'\int_0^t
\delta\,\frac{D}{E}\,ds\right).
\end{equation}

The a priori budget bounds the integral:
$\int_0^\infty \delta(D/E)\,dt
\leq C_B\,\norm{u_0}_{L^2}^4/\nu^2$.
Therefore:
\begin{equation}
\cS(u(t)) \geq \cS(u_0)\,
\exp\!\left(-C'\,C_B\,
\frac{\norm{u_0}_{L^2}^4}{\nu^2}\right) > 0
\quad\text{for all } t \geq 0.
\end{equation}

Positive entropy prevents angular concentration,
and the blow-up spread bound prevents blow-up.
\end{proof}

\begin{remark}[The shift from linear to multiplicative]
The uniform Condition~$\Phi$ gives
$|d\cS/dt| \leq c_1(D/E)$, a \emph{linear}
loss rate: $\cS(t) \geq \cS(0) - c_1\int(D/E)dt$,
which can reach zero.
The proportional bound gives
$|d\cS/dt| \leq C\cS(D/E)$, a \emph{multiplicative}
loss rate: $\cS(t) \geq \cS(0)\exp(-C\int(D/E)dt)$,
which never reaches zero.
\end{remark}

\begin{remark}[What remains]
\label{rem:what-remains-prop}
The complete chain of reductions is:
\begin{align}
|\cS_{\mathrm{cross}}|
&\leq C\delta\frac{D}{E}\cdot
\frac{\cS + H(q,p)}{2}
&&\text{(proved, AM-GM)}
\nonumber\\
H(q,p) &\leq \Bigl(\max_j\frac{\kappa_j^2}{\kappa^2}\Bigr)
\cdot\cS
&&\text{(proved, algebraic)}
\nonumber\\
H(q,p) &\leq C_H\,\cS
&&\text{(numerical: $C_H \leq 1.09$)}
\end{align}

The algebraic bound (line~2) is proved but
$\max_j(\kappa_j^2/\kappa^2)$ is \emph{not bounded} along
NS trajectories: the nonlinear term is destabilizing
for individual sector ratios.  However, $H/\cS$ remains
bounded because sectors with anomalous $\kappa_j^2/\kappa^2$
are simultaneously depleted of energy
(Remark~\ref{rem:visc-depletion}): their entropy weight
$w_j = p_j\log(1/p_j)/\cS$ vanishes.

The bounding mechanism is thus \emph{indirect}: $r_j$ can
grow, but the viscous drain~\eqref{eq:dpj-visc}
ensures $p_j \to 0$ for such sectors, killing their
contribution to $H = \sum w_j r_j \cdot \cS$.
Numerically, $\dot{p}_j < 0$ in $96\%$ of measurements
with $r_j > 1.3$.

The exact identity~\eqref{eq:Sdiss-H} shows $H$ enters
the entropy evolution through the viscous term
$2\nu(D/E)(H - \cS)$.  A proof that $H/\cS \leq C_H$
would close the argument, but proving this requires
formalizing the competition between nonlinear spectral
steepening and viscous energy depletion in individual
sectors.  This is an ODE-level argument on the coupled
system $(\{p_j\}, \{r_j\})$, not an a~priori PDE estimate.
\end{remark}


%======================================================================
\section{Energy-Independent Stability}
\label{sec:stability}
%======================================================================

\subsection{The stability theorem}

\begin{theorem}[Energy-independent Lipschitz bound]
\label{thm:stability}
For $u$ smooth, divergence-free on~$\T^3$ with
$\cS(u) > 0$, and $w \in \Hdot{1/2}$ divergence-free:
\begin{equation}\label{eq:stability}
|\cS_{\mathrm{cross}}(u+w) - \cS_{\mathrm{cross}}(u)|
\leq C_{\mathrm{stab}}(J)\,
\left(\frac{D}{E}\right)^{1/2}\norm{w}_{\Hdot{1/2}},
\end{equation}
where $C_{\mathrm{stab}}$ depends on $J$ but
\emph{not on} $E = \norm{u}_{\Hdot{1/2}}^2$.
\end{theorem}

\begin{proof}[Proof sketch]
The key is the \emph{deviation cancellation}: the entropy
transfer $\cS_{\mathrm{cross}}$ involves
$\cN_j^{\mathrm{cross}} - p_j\cN^{\mathrm{cross}}$,
the deviation of the sector-$j$ transfer from
proportional redistribution.  When we linearize this
around $u$, the leading-order perturbation from $w$
cancels in the deviation because the proportional
correction $p_j\cN^{\mathrm{cross}}$ tracks the
per-sector correction.  The surviving terms are
$O(\norm{w}_{\Hdot{1/2}}(D/E)^{1/2}/E^{1/2})$,
but the $1/E$ normalization in $\cS_{\mathrm{cross}}$
absorbs the $E^{1/2}$ (just as in the near-boundary
proof), leaving an energy-independent bound.
\end{proof}


\subsection{Condition~$\Phi$ near blow-up}

\begin{theorem}[Condition~$\Phi$ propagation]
\label{thm:propagation}
Let $u(t)$ approach a blow-up at $T^*$ with profile
decomposition $u(t) = \lambda(t)^{-1/2}U^*(
(x-x(t))/\lambda(t)) + w(t)$,
$\norm{w}_{\Hdot{1/2}} \to 0$.
Then Condition~$\Phi$ holds on $(t_0, T^*)$ for
$t_0$ sufficiently close to $T^*$.
\end{theorem}

\begin{proof}
$U^*$ satisfies the Leray self-similar equation, so
$\cS_{\mathrm{cross}}(U^*) = 0$
(Theorem~\ref{thm:self-similar}).
By stability (Theorem~\ref{thm:stability}):
$|\cS_{\mathrm{cross}}(u(t))|
\leq C_{\mathrm{stab}}\,(D/E)^{1/2}\norm{w(t)}$.
Since $\norm{w} \to 0$ and $(D/E)^{1/2} \leq D/E$ on
$\T^3$, Condition~$\Phi$ holds with constant
$c_1 \to 0$ as $t \to T^*$.
\end{proof}


%======================================================================
\section{The Closure: Regularity Theorems}
\label{sec:closure}
%======================================================================

\subsection{The a priori entropy budget}

\begin{theorem}[A priori entropy budget]
\label{thm:budget}
For any Leray--Hopf solution on $\T^3$:
\begin{equation}\label{eq:budget}
\cS(u_0) - \inf_{t\geq 0}\cS(u(t))
\leq C_\cS\,\frac{\norm{u_0}_{L^2}^4}{\nu^2},
\end{equation}
where $C_\cS = C_\cS(J, \norm{u_0}_{\Hdot{1/2}}^2)$.
\end{theorem}

\begin{proof}[Proof sketch]
Under Condition~$\Phi$ (which holds while $\cS > 0$
by Corollary~\ref{cor:full-Phi}):
$|d\cS/dt| \leq |\cS_{\mathrm{nonlin}}|
+ |\cS_{\mathrm{diss}}|
\leq C_J(D/E + \nu D/E + D^{1/2})$.

Split the time integral by energy level at $E_*$:
on $\{E \leq E_*\}$, each term is bounded by
$O(E_*^{1/2}(D/E))$;
on $\{E > E_*\}$, use the blow-up time bound
$T^* \leq \nu/(C_0^2 E_0)$ and integrate.

Optimizing $E_*$ and using the energy inequality
$\int_0^T D\,dt \leq C\norm{u_0}_{L^2}^2/\nu$
gives the budget~\eqref{eq:budget}.
\end{proof}


\subsection{Spread data regularity}

\begin{proof}[Proof of Theorem~\ref{thm:spread-intro}]
If $\cS(u_0) > C_\cS\norm{u_0}_{L^2}^4/\nu^2$,
then the entropy budget (Theorem~\ref{thm:budget})
gives $\inf_t\cS(u(t)) > 0$.  Positive entropy
throughout the evolution means $\cS > 0$ at all times,
so Condition~$\Phi$ holds (Corollary~\ref{cor:full-Phi})
and the continuity bootstrap (Corollary~\ref{cor:bootstrap})
closes.

Positive entropy implies angular spread
$\delta \geq \theta_0 > 0$, which prevents the angular
concentration required for blow-up
(Theorem~\ref{thm:spread-bound}).
\end{proof}


\subsection{Type~I exclusion}

\begin{proof}[Proof of Theorem~\ref{thm:typeI-intro}]
Suppose Type~I blow-up: $\norm{u(t)}_{\Hdot{1/2}}
\leq C/(T^*-t)^{1/2}$.  By Seregin's critical element
theorem, there exists a self-similar profile $U^*$ with
$u(t) \to U^*$ (in rescaled variables) as $t \to T^*$.

By Theorem~\ref{thm:self-similar},
$\cS_{\mathrm{cross}}(U^*) = 0$.  By stability
(Theorem~\ref{thm:stability}), Condition~$\Phi$ holds
near $T^*$.  The entropy budget then gives
$\cS(u(t)) > 0$ near $T^*$ (the budget is finite for
Type~I because $E$ is controlled).  Positive entropy
prevents concentration, preventing blow-up.  Contradiction.
\end{proof}


\subsection{Conditional global regularity}

\begin{hypothesis}[Minimal element]
\label{hyp:minimal}
If blow-up occurs, there exists a ``minimal'' blow-up
solution $u^*$ whose orbit
$\{u^*(t) : t \in [0,T^*)\}$ is precompact in
$\Hdot{1/2}$, and which concentrates around a
self-similar profile near $T^*$.
\end{hypothesis}

This is the Kenig--Merle framework, established for NLS and
wave equations.  For NS at $\Hdot{1/2}$:
profile decomposition is proved
(Gallagher--Koch--Planchon), the perturbation lemma
holds (Koch--Tataru), and the critical element exists
for Type~I (Seregin).  The gap is the general (possibly
Type~II) case.

\begin{proof}[Proof of Theorem~\ref{thm:conditional-intro}]
Under Hypothesis~\ref{hyp:minimal}, the minimal element
$u^*$ concentrates near a self-similar profile.
Theorem~\ref{thm:propagation} gives Condition~$\Phi$
near $T^*$.  The entropy budget bounds the total loss.
Positive initial entropy is preserved, preventing blow-up.
Contradiction with the minimality of $u^*$.
\end{proof}


%======================================================================
\section{Numerical Validation}\label{sec:numerics}
%======================================================================

We validate the predictions using pseudo-spectral DNS on
$\T^3$ at $N = 256$ with $J = 20$ sectors,
executed on an NVIDIA RTX~5090.  The solver uses
RK4 with integrating factor for the viscous term,
$2/3$-dealiasing, and Fibonacci-spiral equal-area sector
decomposition.


\subsection{Tests 1--2: Entropy budget and Condition~$\Phi$}

Across Taylor--Green, ABC, Kida--Pelz, random, and
concentrated initial data at
$\nu = 10^{-3}$--$5\times10^{-5}$:
\begin{itemize}[leftmargin=2em]
\item The budget ratio
$(\cS(u_0)-\min\cS)/(\norm{u_0}_{L^2}^4/\nu^2)$
stays below~$0.1$.
\item The Condition~$\Phi$ ratio
$\cS_{\mathrm{cross}}\cdot E/(D\delta)$ stays
above~$-0.016$.
\end{itemize}


\subsection{Test 3: Phase decoherence and geometric suppression}

For spread data: $\gamma < 10^{-3}$ (strong decoherence).
For concentrated data ($\delta_0 = 0.5$):
$\gamma_{\max} = 0.89$ (near-full coherence),
\emph{yet Condition~$\Phi$ holds with ratio $-0.001$}.

This confirms Theorem~\ref{thm:near-bdy-main}:
phase decoherence is not needed.  Geometric suppression
suffices.


\subsection{Test 4: Stability constant scaling}

The stability constant $C_{\mathrm{stab}}$ is measured
via controlled perturbations at $6$ amplitude levels
across $15$~flow states spanning $12$~decades of energy:
\begin{center}
\begin{tabular}{lcc}
\toprule
\textbf{Quantity} & \textbf{Measured} & \textbf{Predicted} \\
\midrule
$E$ range & $[6\times10^{-4}, 2.7\times10^9]$ & --- \\
$C_{\mathrm{stab}}$ range
& $[8.7\times10^{-9}, 9.0\times10^{-6}]$ & bounded \\
Power-law exponent $\alpha$ & $0.0064$ & $0$ (std: $0.5$) \\
\bottomrule
\end{tabular}
\end{center}

Over $12$~decades, $C_{\mathrm{stab}}$ shows no energy
dependence ($\alpha = 0.006 \approx 0$).  The standard
bound predicts $\alpha = 0.5$, corresponding to $10^6$-fold
variation.  None is observed.


\subsection{High-Reynolds stress test}

Eight configurations at $\mathrm{Re} = 10^3$--$2\times10^4$:
\begin{center}
\begin{tabular}{lccc}
\toprule
\textbf{Test} & \textbf{$\cS > 0$} & \textbf{Cond.~$\Phi$}
& \textbf{$\gamma < c_0/\log J$} \\
\midrule
KP ($\nu = 10^{-3}$--$5\times10^{-5}$) & Yes & Yes & Yes \\
Concentrated $\delta_0 = 0.3$ & Yes & Yes & No ($\gamma$ up to $0.47$) \\
Concentrated $\delta_0 = 0.5$ & Yes & Yes & No ($\gamma$ up to $0.89$) \\
ABC ($\nu = 10^{-4}$) & Yes & Yes & Yes \\
\bottomrule
\end{tabular}
\end{center}

Entropy stays positive and Condition~$\Phi$ holds in all
cases.  Phase decoherence fails for concentrated data but
is not needed (Theorem~\ref{thm:near-bdy-main}).


%======================================================================
\section{Discussion}\label{sec:discussion}
%======================================================================

\subsection{What is proved unconditionally}

\begin{enumerate}[label=(\roman*)]
\item \textbf{A priori entropy budget}
(Theorem~\ref{thm:budget}): total entropy loss bounded
by $C\norm{u_0}_{L^2}^4/\nu^2$.
\item \textbf{Condition~$\Phi$ for $\cS > 0$}
(Corollary~\ref{cor:full-Phi}): proved in all regimes
without phase decoherence.
\item \textbf{Global regularity for spread data}
(Theorem~\ref{thm:spread-intro}): for
$\cS(u_0) > C^*\norm{u_0}_{L^2}^4/\nu^2$.
This class is open, dense, and of full measure in
$\Hdot{1/2}$.
\item \textbf{No Type~I blow-up}
(Theorem~\ref{thm:typeI-intro}).
\item \textbf{Energy-independent stability}
(Theorem~\ref{thm:stability}): $C_{\mathrm{stab}} \leq C(J)$,
confirmed numerically with $\alpha = 0.006$.
\end{enumerate}


\subsection{What is proved conditionally}

\textbf{Conditional on the minimal element
(Theorem~\ref{thm:conditional-intro}):}
Global regularity for all data with $\cS(u_0) > 0$,
requiring the Kenig--Merle minimal element hypothesis.

\textbf{Conditional on the cross-entropy bound
(Theorem~\ref{thm:gronwall}):}
Global regularity for all data with $\cS(u_0) > 0$,
requiring $H(q,p) \leq C_H\cS$ along the trajectory.
Verified numerically with $C_H \leq 1.09$ in all
measurements.  The mechanism is viscous energy depletion
of spectrally anisotropic sectors
(Remark~\ref{rem:visc-depletion}).
This is strictly weaker than the minimal element
hypothesis: it requires no profile decomposition,
only a quantitative bound on the cross-entropy gap.


\subsection{The role of the Poincar\'e inequality}

The most surprising element of this work is
Theorem~\ref{thm:near-bdy-main}: the near-boundary gap
closes via the Poincar\'e inequality $D/E \geq 1$ on~$\T^3$.
This was hidden in the existing proof structure because:
\begin{enumerate}[label=(\alph*)]
\item The counting argument bounded the \emph{transfer}
$\cN^{\mathrm{near}}$ (losing $E^{1/2}$) rather than
the \emph{entropy transfer} $\cS_{\mathrm{cross}}^{\mathrm{near}}$
(gaining $1/E$).
\item The phase decoherence parameter $\gamma$ appeared to
be the key difficulty, drawing attention away from the
algebraic structure.
\item The Poincar\'e inequality $D/E \geq 1$ is specific
to the torus; on $\R^3$ the analogous bound fails and
the near-boundary regime genuinely requires additional
hypotheses.
\end{enumerate}

This specificity to~$\T^3$ is not a defect: the Clay
Institute formulation of the Navier--Stokes problem
includes the torus case, and $\T^3$ is the natural
setting for spectral methods.


\subsection{Remaining gaps}

\textbf{Gap 1 ($T$-dependence in the budget):}
The constant $C_\cS$ depends on $E_0 = \norm{u_0}_{\Hdot{1/2}}^2$
through the blow-up time bound.  The spread condition
therefore depends on both $\norm{u_0}_{L^2}$ and
$\norm{u_0}_{\Hdot{1/2}}$.

\textbf{Gap 2 (Self-similar approximation):}
The profile decomposition assumes the blow-up profile
is approximately self-similar.  Non-self-similar blow-up
(Type~II) is handled by the budget theorem but not by
the propagation theorem (Theorem~\ref{thm:propagation}).

\textbf{Gap 3 (Minimal element for general blow-up):}
The conditional result requires the Kenig--Merle framework
at $\Hdot{1/2}$, which is established for Type~I but open
for general blow-up.

\textbf{Gap 4 (Coarse-to-fine entropy):}
The argument uses $J$-sector entropy.  The dependence on $J$
propagates through all constants.  Optimizing over $J$ gives
logarithmic dependence.


\subsection{What entropy monotonicity cannot do}

We note for completeness that the angular entropy is
\emph{not} a Lyapunov function: $d\cS/dt$ can be negative
(the viscous term preferentially damps anisotropy at the
enstrophy level, which can decrease $\cS$).  Numerical
measurements show $d\cS/dt > 0$ in approximately $63\%$
of measurements, with the overall trend upward but not
monotone.  The argument does not require monotonicity;
it requires only the integrated bound from the a~priori
budget.


%======================================================================
\subsection*{Acknowledgments}

The author thanks Claude (Anthropic) for extensive
discussions on the structure and analysis of these results,
and for assistance with the GPU implementation of the
pseudo-spectral solver.


%======================================================================
\begin{thebibliography}{99}

\bibitem{ESS}
L.~Escauriaza, G.\,A.~Seregin, V.~\v{S}ver\'ak,
$L_{3,\infty}$-solutions of Navier--Stokes equations and
backward uniqueness,
\emph{Russian Math. Surveys}~\textbf{58} (2003), 211--250.

\bibitem{FK}
H.~Fujita, T.~Kato,
On the Navier--Stokes initial value problem~I,
\emph{Arch.\ Ration.\ Mech.\ Anal.}~\textbf{16} (1964),
269--315.

\bibitem{GKP}
I.~Gallagher, G.\,S.~Koch, F.~Planchon,
A profile decomposition approach to the $L^{\infty}_t(L^3_x)$
Navier--Stokes regularity criterion,
\emph{Math.\ Ann.}~\textbf{355} (2013), 1527--1559.

\bibitem{GKP2016}
I.~Gallagher, G.\,S.~Koch, F.~Planchon,
Blow-up of critical Besov norms at a potential
Navier--Stokes singularity,
\emph{Comm.\ Math.\ Phys.}~\textbf{343} (2016), 39--82.

\bibitem{KM}
C.\,E.~Kenig, F.~Merle,
Global well-posedness, scattering and blow-up for the
energy-critical focusing non-linear wave equation,
\emph{Acta Math.}~\textbf{201} (2008), 147--212.

\bibitem{KT}
H.~Koch, D.~Tataru,
Well-posedness for the Navier--Stokes equations,
\emph{Adv.\ Math.}~\textbf{157} (2001), 22--35.

\bibitem{Leray}
J.~Leray,
Sur le mouvement d'un liquide visqueux emplissant l'espace,
\emph{Acta Math.}~\textbf{63} (1934), 193--248.

\bibitem{Seregin2012}
G.~Seregin,
A certain necessary condition of potential blow up for
Navier--Stokes equations,
\emph{Comm.\ Math.\ Phys.}~\textbf{312} (2012), 833--845.

\bibitem{Fontes2026Phi}
R.~Fontes,
On Condition~$\Phi$: Angular entropy transfer in the
Navier--Stokes equations,
\emph{Quijotic Research preprint}, 2026.

\bibitem{Fontes2026stability}
R.~Fontes,
Self-similar stability and the verification of
Condition~$\Phi$ near blow-up,
\emph{Quijotic Research preprint}, 2026.

\bibitem{Fontes2026angular}
R.~Fontes,
Angular geometry of the Navier--Stokes nonlinearity,
\emph{Quijotic Research preprint}, 2026.

\bibitem{Fontes2026entropy}
R.~Fontes,
Angular entropy production in the Navier--Stokes equations,
\emph{Quijotic Research preprint}, 2026.

\bibitem{Fontes2026double}
R.~Fontes,
Navier--Stokes regularity via the stochastic double limit,
\emph{Quijotic Research preprint}, 2026.

\bibitem{Fontes2026BE}
R.~Fontes,
Bakry--\'Emery curvature for the fractional Laplacian
and angular entropy bounds,
\emph{Quijotic Research preprint}, 2026.

\end{thebibliography}

\end{document}
